\begin{remark}
  From now on, assume basic arithmetic/algebraic properties of $\R$.
\end{remark}

\subsection{Sequences}

\begin{definition}{Sequence}{sequence}
  A \textbf{sequence} is a function $S:\N\to\R$.
  \begin{flushleft}
    Notation: $S(n)$ or $S_n$\\
    or $(S_n)_{n=1}^{\infty}$ or $(S_n)_{n\in\N}$.
  \end{flushleft}
\end{definition}

\begin{example}
  \begin{gather*}
    S=\left(\frac1n\right)_{n\in\N}
    =\left(1,\frac12,\frac13,\ldots\right),\\
    S=\bigl((-1)^n\bigr)_{n\in\N}=(-1,1,-1,1,\ldots).
  \end{gather*}
\end{example}

\begin{note}
  Notation: $(\,)$ sequences, $\{\,\}$ sets.\\
  $(1,-1)$ is not $\{1,-1\}$.
\end{note}

\subsection{Convergence}

\begin{definition}{Convergence}{sequence-convergence}
  \textbf{Important!}
  \begin{flushleft}
    Let $S:\N\to\R$ be a sequence,\\
    let $t\in\R$.\\[0.5em]
    $S$ \textbf{converges to $t$}, or $S\to t$,
    or $\lim_{n\to\infty}S_n=t$,
  \end{flushleft}
  means
  \[
    \begin{gathered}
      \forall\varepsilon\in\R,\ \varepsilon>0,\\[0.5em]
      \exists N\in\N\text{ s.t.}\\[0.5em]
      n>N\implies\abs{S_n-t}<\varepsilon.
    \end{gathered}
  \]
\end{definition}

\begin{intuition}
  \begin{flushleft}
    $S_n\to t\iff$\\[0.5em]
    for every $\varepsilon>0$ (small window) challenging our assertion,\\[0.5em]
    we can find a $N\in\N$ (concrete index)\\[0.5em]
    such that for every $n>N$, $S_n$
    has sequence output: $\abs{S_n-t}<\varepsilon$.\\[0.5em]
    Distance $S_n\leftrightarrow t$
    is smaller than $\varepsilon$ window.
  \end{flushleft}
  \begin{center}
    \begin{tikzpicture}
      \draw[number line] (-3,0) -- (3,0);
      \draw[sketch accent] (-1.8,0) -- (1.8,0);
      \foreach \x/\label in {-1.8/{t-\varepsilon},1.8/{t+\varepsilon}} {
        \draw[orange!75!black,thick,fill=orange!5] (\x,0) circle (2.5pt);
        \node[above=6pt] at (\x,0) {$\label$};
      }
      \draw (0,-0.1) -- (0,0.1) node[below=5pt] {$t$};
      \node[sketch point] at (0.8,0) {};
      \node[below] (sn) at (1.4,-0.65) {$S_n$};
      \draw[->,>=latex] (sn) -- (0.8,-0.1);
    \end{tikzpicture}
  \end{center}
\end{intuition}

\begin{example}[$1/n\to0$]
  For $S=(1/n)_{n\in\N}$:
  \begin{center}
    \begin{tikzpicture}[x=6cm]
      \draw[number line] (-0.15,0) -- (1.2,0);
      \draw (0,-0.1) -- (0,0.1) node[below=5pt] {$0$};
      \foreach \n in {1,...,12} {
        \draw ({1/\n},-0.08) -- ({1/\n},0.12);
      }
      \foreach \n in {1,2,3,4} {
        \node[below=5pt] at ({1/\n},0) {$\frac1{\n}$};
      }
      \node[above] at (0.04,0.05) {$\cdots$};
      \draw[->,>=latex,sketch accent] (0.4,0.5) -- (0.08,0.5);
    \end{tikzpicture}
  \end{center}

  $S_n\to0$: looks like so on the graph.

  \textbf{Sketch (intuition).} WTS
  (Definition~\ref{def:sequence-convergence}):
  \[
    \begin{gathered}
      \forall\varepsilon>0,\quad\exists N\in\N,\\[0.5em]
      \forall n>N,\quad\left|\frac1n-0\right|<\varepsilon.
    \end{gathered}
  \]
  \begin{flushleft}
    Want to construct $N$ to
    answer challenging $\varepsilon$\\
    by showing $\abs{1/n}<\varepsilon$ for $n>N$.
  \end{flushleft}

  \begin{proof}
    $\forall\varepsilon>0$:\\
    Choose $N\in\N$ s.t.\ $1/N<\varepsilon$.
    Archimedean property (Theorem~\ref{thm:archimedean-property}).\\
    Also, if $n>N>0$, then $1/n<1/N$.
    \[
      \left|\frac1n-0\right|=\frac1n<\frac1N<\varepsilon.\qedhere
    \]
  \end{proof}
\end{example}
